\section{Indexed Logits CUDA Implementation}
\label{app:IndexedLogits}

This appendix describes the indexed logits CUDA extension used in our experiments and reports benchmark results for the implementation currently released in the accompanying repository.\footnote{\url{https://github.com/pettyjohnjn/indexed_logits}} The key operation is
\begin{equation}
\mathrm{out}[i,j] = \langle H[i,:], W[\mathrm{idx}[i,j],:] \rangle,
\end{equation}
where $H \in \mathbb{R}^{N \times d}$ is a matrix of hidden states, $W \in \mathbb{R}^{V \times d}$ is the vocabulary projection matrix, and $\mathrm{idx} \in \mathbb{Z}^{N \times k}$ selects a subset of vocabulary rows for each position.

\subsection{Implementation}

The released implementation is a correctness-first fused CUDA kernel. It avoids materializing the tensor $W[\mathrm{idx}] \in \mathbb{R}^{N \times k \times d}$ and instead computes each output element directly in a single thread:
\begin{equation}
\mathrm{out}[i,j] = \sum_{t=1}^{d} H[i,t] \, W[\mathrm{idx}[i,j], t].
\end{equation}
Accumulation is performed in FP32 even when the inputs are FP16 or BF16, and the result is cast back to the input dtype when written to the output tensor.

The backward pass computes two gradients:
\begin{align}
\nabla_H[i,t] &= \sum_{j=1}^{k} \nabla_{\mathrm{out}}[i,j] \, W[\mathrm{idx}[i,j], t], \\
\nabla_W[v,t] &= \sum_{(i,j): \mathrm{idx}[i,j]=v} \nabla_{\mathrm{out}}[i,j] \, H[i,t].
\end{align}
The implementation computes $\nabla_H$ with a separate CUDA kernel and accumulates $\nabla_W$ into an FP32 buffer using atomic additions. This design is simple and robust under mixed precision, though it is not a heavily tuned GEMM replacement.

\subsection{Scope of the Current Repository}

The current public repository supports:
\begin{itemize}
\item fused indexed-logit forward computation;
\item custom backward kernels for $\nabla_H$ and $\nabla_W$;
\item FP32 accumulation in forward and weight-gradient computation;
\item automatic differentiation through a PyTorch custom autograd function.
\end{itemize}

It does \emph{not} currently include a LoRA-specific fused kernel, a Triton fallback, or the more aggressive shared-memory and tiling optimizations that would be needed to justify stronger kernel-efficiency claims. Accordingly, the experimental claims below are limited to memory efficiency, correctness, and end-to-end performance relative to straightforward PyTorch baselines.

\subsection{Correctness}

We evaluated the CUDA extension against a PyTorch reference implementation that explicitly materializes $W[\mathrm{idx}]$ and uses standard autograd for the backward pass. Experiments were run on a single A100-40GB GPU with PyTorch 2.8.0 and CUDA 12.8/12.9 runtime compatibility.

Across representative FP16 configurations, forward outputs and gradients matched the PyTorch reference closely. For example, in the extended FP16 sweep:
\begin{itemize}
\item at $(N,d,k)=(4096,768,16)$, the mean absolute error was $4.6\times 10^{-3}$ for the forward pass, $1.4\times 10^{-10}$ for $\nabla_H$, and $3.5\times 10^{-5}$ for $\nabla_W$;
\item at $(4096,768,128)$, the mean absolute error was $4.6\times 10^{-3}$ for the forward pass, $2.8\times 10^{-8}$ for $\nabla_H$, and $1.1\times 10^{-3}$ for $\nabla_W$;
\item at $(4096,768,512)$, the mean absolute error was $4.6\times 10^{-3}$ for the forward pass, $3.5\times 10^{-7}$ for $\nabla_H$, and $4.4\times 10^{-3}$ for $\nabla_W$.
\end{itemize}

BF16 results showed the expected increase in forward quantization error while remaining numerically stable. For example, in the main BF16 benchmark suite the forward mean absolute error was approximately $3.6\times 10^{-2}$ for GPT-2-like configurations with $d=768$, and $5.2\times 10^{-2}$ for the $d=1536$ configuration. Weight-gradient mean absolute error remained below $1.8\times 10^{-2}$ in all BF16 cases we measured.

\subsection{Benchmark Methodology}

We compare the fused kernel against two baselines:
\begin{enumerate}
\item \textbf{Naive subset}: materialize $W[\mathrm{idx}] \in \mathbb{R}^{N \times k \times d}$ and compute logits with elementwise multiplication and reduction;
\item \textbf{Dense GEMM}: compute the full logits matrix $H W^\top \in \mathbb{R}^{N \times V}$ and then gather the requested subset.
\end{enumerate}

All timings were collected with CUDA events after warmup iterations. Peak memory is reported using PyTorch CUDA memory statistics. We report forward-only latency and end-to-end forward+backward latency. For the extended sweep we additionally ran three random seeds and report consistent trends across seeds.

\subsection{Main Results}

Table~\ref{tab:indexed_logits_main} summarizes representative FP16 results on an A100-40GB GPU. The most consistent pattern is that the fused kernel is substantially more memory-efficient than the naive subset baseline and provides a reliable end-to-end speedup over that baseline. Relative to dense GEMM, the fused kernel is usually faster in the regimes we care about, though the margin is modest when $k$ becomes large.

\begin{table}[t]
\centering
\caption{Indexed logits benchmark results on a single A100-40GB GPU (FP16).}
\label{tab:indexed_logits_main}
\small
\begin{tabular}{lrrrrr}
\toprule
Config $(N,d,k)$ & Method & Forward (ms) & Fwd+Bwd (ms) & Peak Mem (MB) & Relative Note \\
\midrule
$(4096,768,32)$
  & Dense GEMM   & 4.81 & 12.89 & 926.4 & Computes full $V$ logits \\
  & Naive subset & 0.89 & 2.53  & 696.3 & Materializes $N \times k \times d$ \\
  & Fused kernel & 1.02 & 1.84  & 324.0 & 1.38$\times$ faster than naive end-to-end \\
\midrule
$(4096,768,128)$
  & Dense GEMM   & 4.84 & 12.98 & 932.3 & Full-vocabulary baseline \\
  & Naive subset & 3.50 & 9.39  & 2531.1 & High intermediate-memory cost \\
  & Fused kernel & 3.34 & 6.46  & 342.6 & 1.45$\times$ faster, 7.39$\times$ less memory than naive \\
\midrule
$(4096,768,256)$
  & Dense GEMM   & 4.88 & 13.12 & 940.7 & Full-vocabulary baseline \\
  & Naive subset & 6.97 & 18.40 & 4955.4 & Very large intermediate tensor \\
  & Fused kernel & 6.76 & 12.59 & 346.8 & 1.46$\times$ faster, 14.3$\times$ less memory than naive \\
\midrule
$(4096,1536,128)$
  & Dense GEMM   & 9.55 & 25.48 & 1015.8 & Full-vocabulary baseline \\
  & Naive subset & 7.05 & 18.83 & 5036.8 & Intermediate dominated \\
  & Fused kernel & 6.96 & 13.33 & 664.0 & 1.41$\times$ faster than naive end-to-end \\
\midrule
$(8192,768,128)$
  & Dense GEMM   & 9.77 & 25.92 & 1772.0 & Full-vocabulary baseline \\
  & Naive subset & 6.97 & 18.44 & 4968.0 & Intermediate dominated \\
  & Fused kernel & 6.76 & 12.71 & 359.4 & 1.45$\times$ faster, 13.8$\times$ less memory than naive \\
\bottomrule
\end{tabular}
\end{table}

These results support three conclusions.

\paragraph{1. The main win is memory efficiency.}
The fused kernel consistently avoids the large $N \times k \times d$ intermediate tensor that dominates the naive subset implementation. The advantage grows quickly with $k$. In the extended $k$ sweep at $(N,d,V)=(4096,768,50257)$, peak memory for the fused kernel rose only from 333.8 MB at $k=16$ to 355.2 MB at $k=512$, while the naive subset baseline grew from 404.3 MB to 9816.6 MB.

\paragraph{2. End-to-end speedup over naive subset is robust.}
The fused kernel provides a consistent end-to-end speedup over the naive subset baseline across all measured configurations. In the main FP16 suite, the forward+backward speedup over naive subset ranged from 1.38$\times$ to 1.46$\times$.

\paragraph{3. Dense GEMM remains a strong baseline.}
The fused kernel is not universally faster than dense GEMM in every forward-only regime. For example, at $(4096,768,256)$ dense GEMM forward time was 4.88 ms versus 6.76 ms for the fused kernel. However, dense GEMM pays a fixed cost for computing all $V$ logits and uses substantially more memory. In the forward+backward measurements used most directly for training-time analysis, the fused kernel was competitive with or faster than dense GEMM in every configuration we evaluated.

\subsection{Scaling with Subset Size}

The extended sweep over subset size $k$ shows the expected scaling behavior. At fixed $(N,d,V)=(4096,768,50257)$:
\begin{itemize}
\item fused forward time increased from 0.53 ms at $k=16$ to 13.57 ms at $k=512$;
\item fused forward+backward time increased from 1.20 ms at $k=16$ to 24.76 ms at $k=512$;
\item dense GEMM remained nearly flat, with forward time around 4.8--5.0 ms and forward+backward time around 12.9--13.4 ms;
\item naive subset scaled roughly linearly in both runtime and memory, reaching 9.8 GB peak memory at $k=512$.
\end{itemize}

This scaling trend is exactly the practical setting where indexed logits are useful: when only a moderate subset of logits is needed per position, the fused kernel avoids the large intermediate tensors that make the naive subset formulation unattractive.

\subsection{Repeated Indices and Atomic Contention}

The backward pass accumulates $\nabla_W$ using atomic additions, so one might expect repeated indices to hurt performance. Empirically, however, the opposite trend appeared in our measurements for the current implementation. In a dedicated collision study with $(N,d,V,k)=(4096,1024,50257,64)$ and three seeds:
\begin{itemize}
\item fused forward time decreased from 2.290 ms when indices were sampled from the full vocabulary to 1.604 ms when indices were sampled from only 64 vocabulary rows;
\item fused backward-only time decreased from 2.078 ms to 1.781 ms over the same sweep;
\item fused forward+backward time decreased from 4.617 ms to 3.635 ms.
\end{itemize}

The most plausible explanation is that repeated indices improve reuse of the corresponding rows of $W$, and in this implementation the locality benefit is larger than the additional atomic contention cost. We therefore do \emph{not} claim that duplicate indices are negligible because they are rare; instead, the behavior of repeated indices should be treated as an empirical property of the kernel and hardware combination.

\subsection{PyTorch Integration}

The extension is exposed through a PyTorch custom autograd function and can be used as a drop-in replacement for subset logit computation inside training or evaluation loops:
\begin{verbatim}
out = indexed_logits(H, W, idx)
\end{verbatim}
The interface supports CUDA tensors in FP16 and BF16, accepts integer index tensors, and returns gradients for $H$ and $W$ automatically through PyTorch autograd.

\subsection{Summary}

For the current implementation, the strongest supported claim is not that indexed logits outperforms all dense alternatives in every regime. Rather, the supported claim is that a simple fused CUDA kernel can eliminate the $N \times k \times d$ intermediate tensor, substantially reduce memory usage, and deliver a reliable end-to-end speedup over naive subset computation while remaining numerically accurate under mixed precision.
